\begin{abstract}
We characterize the finite-sample minimax excess variance of Horvitz--Thompson estimation under covariate-dependent experimental assignment. Covariates are independent uniform draws on the unit cube, and the conditional potential-outcome midpoint mean is centered, has an exact main-and-pair-effect decomposition, and satisfies a pooled nonperiodic Sobolev restriction-norm budget. For every sample size \(n\geq2\), covariate dimension \(d\geq2\), and smoothness \(0<s\leq1\), the minimax scaled excess variance is comparable to \(\min\{1,(B/n)^s\}\), where \(B=\binom d2\) is the number of coordinate pairs, with positive comparison constants uniform over this parameter range. A single fair, outcome-independent assignment law, independent of smoothness, attains this scale on all original units. An independent coefficient-sign prior supplies the matching lower bound against every admissible design. The characterization is logarithm-free, includes the endpoint \(s=1\), and yields vanishing minimax scaled excess exactly when \(d=o(\sqrt n)\) for each fixed smoothness. The causal interpretation applies to arbitrary square-integrable potential outcomes with the specified uniform-covariate prognostic class.
\end{abstract}

\section{Introduction}

This paper establishes a finite-sample capacity characterization for covariate balancing in randomized experiments. Under independent uniform covariates, we consider centered prognostic functions consisting exactly of canonical main effects and pairwise interactions, with a common squared nonperiodic Sobolev restriction-norm budget pooled across components. Writing \(n\) for the number of experimental units, \(d\) for covariate dimension, and \(s\) for component smoothness, the best achievable scaled excess variance of the unit-weight Horvitz--Thompson estimator for this class is of order
\[
\min\left\{1,\left(\frac{\binom d2}{n}\right)^s\right\},
\qquad n,d\geq2,\quad 0<s\leq1.
\]
One fair, outcome-independent assignment procedure attains this scale simultaneously across the stated smoothness range. The matching converse applies to every admissible covariate-dependent assignment law.

The statistical question is how much precision assignment can deliver when the outcome regression has many potentially relevant interactions. Pairwise structure associates each component with at most two covariates, while allowing the number of components to grow quadratically with ambient dimension. A pooled regularity budget bounds their total squared component norm and permits its allocation across coordinates and pairs to vary. Our characterization identifies \(B=\binom d2\), the coordinate-pair count, relative to sample size as the quantity governing worst-case excess variance. Below saturation, \(B<n\), the rate is \((B/n)^s\); at and above saturation, it is constant order. For every fixed smoothness, vanishing minimax scaled excess is equivalent to \(d=o(\sqrt n)\), including along arbitrary integer dimension sequences.

The design problem fixes the estimator and specifies both the available information and the population comparison. An admissible design is a Borel assignment kernel using the complete original covariate array, with conditional treatment probability one-half for each unit and assignment independent of potential outcomes given covariates. Every original unit receives a treatment sign and enters the Horvitz--Thompson estimator at its original weight. Performance is evaluated by expected squared prognostic imbalance, averaging over covariate sampling and assignment. The supremum over prognostic functions lies outside that expectation, so each function is fixed across sample realizations. These choices define the minimax signing risk \(R_s(n,d)\), the scaled population imbalance criterion in \cref{obj:def:criterion}, over the function and design classes in \cref{obj:def:sobolev-class,obj:def:design-class}.

The causal interpretation follows from the excess-variance identity of \citet[Proposition 2.3, equation (2.4), under Assumption 2.1 and Definition 2.2]{Cytrynbaum2026Limits}. For independent original units, admissible assignment, and finite potential-outcome second moments, scaled Horvitz--Thompson variance above the law-specific benchmark equals scaled expected squared imbalance of the conditional potential-outcome midpoint mean. Every prognostic function in our class is realized by a Gaussian outcome completion. Consequently, the minimax excess over all uniform-covariate outcome laws with square-integrable potential outcomes and the specified prognosis equals \(R_s(n,d)\), as established in \cref{obj:lem:published-prognostic-capacity,obj:thm:capacity-answer}. The Gaussian subclass realizes the same capacity and provides a concrete calibration with scaled variance benchmark \(4.01\).

The closest rate comparison is \citet[Theorem B.8(ii), Lemma B.9, and Theorem 4.9]{Cytrynbaum2026Limits}. His correctly specified pairwise model admits a structured Gram--Schmidt Walk upper bound of order \((d^2\log n/n)^{s\wedge1}\), while the full-dimensional smooth-class lower bound specializes to exponent \(s\) in two dimensions. For our independent uniform sample and centered pooled order-two Sobolev class, \cref{obj:thm:capacity-answer,obj:thm:log-free-bivariate-capacity} matches upper and lower bounds with explicit interaction-count dependence and a logarithm-free rate. The comparison includes the frequency-rich regime \(B/n\to0\) and the endpoint \(s=1\). An existing random-matrix route supplies one-frequency saturation through \citet[Theorem 1.2(2a)]{KoganNandyHuang2024Extremal}; our characterization covers both that regime and the subcritical frequency-rich comparison at finite sample sizes.

The construction connects minimax experimental design \citep[Sections 2.2--2.4]{Kallus2018} to randomized balancing. In the subcritical regime, independent coordinate reflections and a common random shift supply ordered real spectral features. A finite boundary-rounding procedure protects an initial segment of these features while preserving coordinate means; the protected segment contracts as assignment coordinates reach their boundaries. The feature ordering is independent of smoothness. Combining its row-imbalance guarantees with the reflected nonperiodic Sobolev budget yields simultaneous attainment. In the saturated regime, the procedure uses independent fair signs. The complete assignment law and its guarantees are specified in \cref{obj:def:capacity-handle,obj:def:ordered-boundary-rounding,obj:thm:log-free-bivariate-capacity} under exact-real arithmetic conventions.

The converse uses a finite pair-cosine family with independent coefficient signs drawn before the covariate sample. Every realization satisfies the pooled Sobolev budget. A sample-size-dependent frequency cutoff supplies enough feature directions for a signing lower bound, including directions from pairs sharing an original coordinate. Averaging over this legal prior preserves the fixed-function population comparison and yields the full-design lower bound in \cref{obj:thm:full-design-lower}. Together, the two comparisons have constants bounded uniformly over \(0<s\leq1\). The uniform subclass also transfers the converse and the necessary dimension condition to the larger bounded-density outcome classes specified in \cref{obj:lem:published-prognostic-capacity}.

The main capacity results are stated in \cref{obj:thm:capacity-answer,obj:thm:log-free-bivariate-capacity}.

The paper proceeds through related work, setup and assumptions, main results, the smoothness-independent assignment procedure, and discussion and extensions. The appendices develop causal variance transfer, spectral and rounding bounds, the independent-prior converse, and proof assembly, followed by proofs of the main results.
